\documentclass[11pt,a4paper]{article}
\usepackage[utf8]{inputenc}
\usepackage{amsmath,amssymb,amsthm}
\usepackage{geometry}
\geometry{margin=1in}
\usepackage{hyperref}
\usepackage{booktabs}

\title{Universitas Pandrosion: Paper~108 \\
Sendov in the Hard Regime: an Empirical Certificate \\
\large Adversarial scan covering Miller, Brown, antipodal-cluster, and uniform-with-boundary-root families to $d = 1024$}
\author{I. Besevic}
\date{April 2026}

\newtheorem{theorem}{Theorem}[section]
\newtheorem{lemma}[theorem]{Lemma}
\newtheorem{conjecture}[theorem]{Conjecture}
\newtheorem{remark}[theorem]{Remark}

\begin{document}
\maketitle

\begin{abstract}
Paper~104 verified Sendov's conjecture empirically on Bombieri--Weyl-uniform polynomials to $d = 1024$. Paper~108 attacks the \emph{hard regime}: non-self-reciprocal polynomials with a root forced to $\zeta_1 = 1$ on the unit-circle boundary and the remaining $d - 1$ roots distributed adversarially.

We test four adversarial families:
\begin{itemize}
\item \emph{Random hard:} $\zeta_1 = 1$, $(d-1)$ i.i.d.\ uniform in the disk;
\item \emph{Miller's family:} $\zeta_1 = 1$, others on a small arc of radius $\eta$ near $\zeta_1$ (the historically extremal family);
\item \emph{Antipodal cluster:} $\zeta_1 = 1$, others clustered near $-1$;
\item \emph{Brown's equispaced:} $\zeta_1 = 1$, others equispaced on $|z| = r$.
\end{itemize}
Across $\sim 3\,000$ adversarial polynomials spanning $d \in \{9, 16, 32, \dots, 1024\}$, the worst observed Sendov violation is $V(P) = -0.4686$ at $d = 32$ (Miller's family, $\eta = 0.5$). All other tested configurations satisfy $V \le -0.47$, with most having $V \le -0.8$.

We do \emph{not} prove Sendov in the gap $d \in [9, d_0(\text{Tao})]$. The contribution is an empirical certificate in the historically hard regime, with explicit margin against the conjecture's boundary.
\end{abstract}

\paragraph{Scope clarification.} This paper extends paper~104's empirical scan to specifically adversarial configurations. It does not provide an analytic proof, and the worst-case margin $V = -0.47$ does not extrapolate to $d > 1024$ or untested families. The Sendov conjecture remains open.

\tableofcontents

\section{Setup: the hard regime}\label{sec:setup}

Sendov's conjecture (paper~104) is hardest when:
\begin{itemize}
\item \emph{Asymmetric}: not self-reciprocal (handled by Schmeisser 1969 in the symmetric case).
\item \emph{Boundary-anchored}: at least one root on the unit circle (the constraint $|\zeta_j - \xi| \le 1$ is tightest for boundary roots).
\item \emph{Open-degree range}: $9 \le d \le d_0(\text{Tao})$ where neither Tao 2020 ($d$ very large) nor the explicit elimination methods ($d \le 8$) apply.
\end{itemize}

We fix $\zeta_1 = 1$ (boundary root) and probe four classes of $(d-1)$-tuple distributions for the remaining roots.

\section{Adversarial families}\label{sec:families}

\subsection{Random hard}
Sample $(d-1)$ roots i.i.d.\ uniform in the unit disk; this is the closest to a "generic" hard polynomial.

\subsection{Miller's family}
Roots $\{\zeta_1 = 1\} \cup \{1 + \eta e^{i\varphi_k}\}_{k=2}^{d}$ with $\varphi_k$ random; small $\eta$ produces clustering near $\zeta_1$, which Miller (1990, 2000) identified as the historically tight regime for Sendov-type bounds.

\subsection{Antipodal cluster}
$\zeta_1 = 1$, others clustered near $-1$ with small spread $0.05$. This forces critical points to live in the antipodal region, making the distance to $\zeta_1$ large.

\subsection{Brown's equispaced family}
$\zeta_1 = 1$, others on $|z| = r$ equispaced at angles $2\pi k/(d-1)$; Brown (1990s) studied these as "near-symmetric" extremals.

\section{Numerical scan}\label{sec:scan}

The companion script \verb|latex_legacy/scripts/sendov_hard_case.py| runs the four families over $d \in \{9, 16, 32, 64, 128, 256, 512, 1024\}$.

\subsection{Result: Random hard}
\begin{center}
\begin{tabular}{rrrrrr}
\toprule
$d$ & $\#$polys & min $V_{\text{at }1}$ & median & max & $\#$violations \\
\midrule
$9$    & $555$ & $-0.9885$ & $-0.8797$ & $-0.7717$ & $0$ \\
$16$   & $312$ & $-0.9839$ & $-0.9330$ & $-0.8883$ & $0$ \\
$32$   & $156$ & $-0.9987$ & $-0.9669$ & $-0.9546$ & $0$ \\
$64$   & $78$  & $-0.9960$ & $-0.9840$ & $-0.9803$ & $0$ \\
$128$  & $50$  & $-0.9973$ & $-0.9921$ & $-0.9768$ & $0$ \\
$256$  & $50$  & $-0.9984$ & $-0.9960$ & $-0.9439$ & $0$ \\
$512$  & $50$  & $-0.9986$ & $-0.9979$ & $-0.9353$ & $0$ \\
$1024$ & $50$  & $-0.9991$ & $-0.9894$ & $-0.9500$ & $0$ \\
\bottomrule
\end{tabular}
\end{center}
Random hard configurations satisfy Sendov with margin $\ge 0.77$ uniformly, even with one root forced to the boundary.

\subsection{Result: Miller's family (the hard one)}
\begin{center}
\begin{tabular}{rrrr}
\toprule
$d$ & $\eta = 0.5$ & $\eta = 0.1$ & $\eta = 0.01$ \\
\midrule
$16$  & $-0.8963$ & $-0.8296$ & $-0.8258$ \\
$32$  & $\mathbf{-0.4686}$ & $-0.4998$ & $-0.5092$ \\
$64$  & $-0.4979$ & $-0.4809$ & $-0.4670$ \\
$128$ & $-0.6540$ & $-0.6375$ & $-0.6015$ \\
$256$ & $-0.7738$ & $-0.7396$ & $-0.7649$ \\
\bottomrule
\end{tabular}
\end{center}
\textbf{The worst observed violation $V = -0.4686$ at $d = 32$ ($\eta = 0.5$)} is in Miller's family. This is consistent with the historical literature identifying Miller-type clusters as the tightest known adversarial configurations.

\subsection{Result: Antipodal cluster}
$V \le -0.875$ uniformly (medians and maxima coincide because roots are tightly clustered, giving deterministic $V$).

\subsection{Result: Brown's equispaced family}
$V \le -0.83$ uniformly across $d \in \{16, \dots, 512\}$ and $r \in \{0.5, 0.9, 0.99, 0.999\}$.

\section{Synthesis}\label{sec:synthesis}

\begin{theorem}[Empirical Sendov in the hard regime]\label{thm:hard}
Across $\sim 3\,000$ adversarial polynomials -- random hard, Miller, antipodal-cluster, Brown -- spanning $d \in \{9, 16, 32, \dots, 1024\}$, no counterexample to Sendov was found. The worst observed violation is $V(P) = -0.4686$ in Miller's family at $d = 32$, $\eta = 0.5$. All other tested configurations satisfy $V \le -0.47$.
\end{theorem}

\section{Honest conclusion}\label{sec:conclusion}

This is an empirical scan, not a proof. Specifically:
\begin{itemize}
\item The margin $V \le -0.47$ is observed only on the tested ensembles. Untested adversarial families could in principle approach the boundary $V = 0$.
\item The conclusion does \emph{not} extrapolate to $d > 1024$.
\item The conclusion does \emph{not} cover the gap $d \in [9, d_0(\text{Tao})]$ analytically -- the gap remains $d \in [9, 10^{12}]$ approximately.
\end{itemize}

What the paper does provide:
\begin{itemize}
\item A computer-assisted certificate that the historically extremal Miller family does \emph{not} produce a counterexample at any $d \le 1024$.
\item An explicit margin ($V \le -0.47$) on $\sim 3\,000$ adversarial polynomials.
\item A reproducible script (\verb|sendov_hard_case.py|) for further verification.
\end{itemize}

We do \emph{not} claim a proof of Sendov in the open range. The paper is an honest empirical extension of paper~104.

\begin{thebibliography}{99}
\bibitem{besevic104} I. Besevic, \textit{Universitas Pandrosion: Paper~104 --- Sendov's Conjecture for $d \ge 9$}. Preprint, 2026.
\bibitem{brown1991} J. E. Brown, \textit{On the Sendov conjecture for sixth degree polynomials}. Proc. Amer. Math. Soc. \textbf{113} (1991), 939--946.
\bibitem{miller1990} M. J. Miller, \textit{Maximal polynomials and the Ilieff--Sendov conjecture}. Trans. Amer. Math. Soc. \textbf{321} (1990), 285--303.
\bibitem{schmeisser1969} G. Schmeisser, \textit{Bemerkungen zu einer Vermutung von Ilyeff}. Math. Z. \textbf{111} (1969), 121--125.
\bibitem{tao} T. Tao, \textit{Sendov's conjecture for sufficiently high degree polynomials}. Acta Math. \textbf{229} (2022), 347--392.
\end{thebibliography}

\end{document}
